\section{Introduction}

Evidence-grounded medical question answering requires systems to retrieve relevant biomedical evidence before producing or selecting an answer. Retrieval-augmented generation (RAG) provides a general framework for combining retrieved evidence with language models \citep{lewis2020rag}, and recent medical benchmarking work has shown that retrieval choices strongly affect medical QA performance \citep{xiong2024benchmarking}. However, standard lexical or dense retrieval often ranks passages by pairwise similarity only. Biomedical evidence is frequently organized through higher-order relations among questions, passages, MeSH concepts, biomedical entities, documents, and relation sources.

This paper studies a practical, domain-specific evidence control layer rather than a new generic RAG paradigm. We present KCH-MedRank, a knowledge-constrained local hypergraph learning-to-rank framework that acts as a white-box evidence selection module for biomedical RAG. The method keeps candidate generation reproducible with BM25, dense retrieval, and Reciprocal Rank Fusion \citep{robertson2009bm25,cormack2009rrf}, then builds a local cross-granularity hypergraph around each question and its top candidate passages.

The contribution is a focused and interpretable evidence reranking framework. KCH-MedRank uses biomedical semantic reranking scores, MeSH hierarchy-aware features, local hypergraph diffusion and centrality, and query-group LambdaMART ranking to improve evidence retrieval and expose the evidence signals used by the reranker. Hypergraph signals are used as interpretable ranking features rather than as a hand-tuned standalone scoring formula. The system does not claim to replace clinical judgment, provide clinical diagnosis, or guarantee answer correctness.

The experiments use public biomedical QA resources and emphasize evidence retrieval metrics, structural ablations, and controlled supporting analyses. On the BioASQ evidence retrieval setting, enhanced candidate generation raises the held-out top-100 recall ceiling, and KCH-MedRank improves top-10 evidence retrieval over original Hybrid RRF and a true MedCPT Cross-Encoder baseline. A structural ablation shows that the hypergraph topology alone introduces noise, but medical knowledge constraints reverse this effect into a statistically significant contribution across all top-10 metrics. PubMedQA and BEIR NFCorpus are used as secondary diagnostics: PubMedQA tests evidence coverage in a different QA format, NFCorpus tests retrieval-only robustness on an official BEIR biomedical split, and a fixed-prompt Qwen3-8B pilot is retained as an appendix-style check rather than a main result.

The contributions are:
\begin{itemize}
  \item a reproducible two-stage biomedical evidence retrieval pipeline with BM25, dense retrieval, Hybrid RRF, enhanced candidate fusion, and held-out evaluation;
  \item a local knowledge-constrained hypergraph representation connecting question, passage, document, entity, MeSH concept, hierarchy, and auxiliary relation signals;
  \item a learning-to-rank reranker that combines retrieval, biomedical semantic, MeSH hierarchy, entity, relation, and hypergraph features under query-level groups;
  \item a structural ablation framework showing that medical knowledge constraints are necessary and sufficient for the hypergraph to improve retrieval (significant Knowledge$\times$Structure interaction, $p < 0.01$);
  \item an evidence-oriented evaluation including Recall@k, Hit@k, MRR@k, nDCG@k, hard reranking subsets, paired bootstrap significance tests, failure analysis, external retrieval-only diagnostics, and a controlled PubMedQA Qwen3-8B appendix pilot.
\end{itemize}

The failure analysis also shows lost gold passages, so the paper treats the method as a net retrieval improvement with clear remaining failure modes rather than as a complete solution for clinical medical QA.
